% Generated draft: requires mathematical and source review.
\documentclass[11pt]{article}
\usepackage[margin=25mm]{geometry}
\usepackage{fontspec}
\setmainfont{Latin Modern Roman}
\usepackage{amsmath,amssymb,mathtools,unicode-math}
\setmathfont{Latin Modern Math}
\usepackage{enumitem,needspace,xurl,hyperref}
\setlength{\parindent}{0pt}
\setlength{\parskip}{5pt}
\setlength{\emergencystretch}{3em}
\newcommand{\N}{\mathbb{N}}
\newcommand{\Z}{\mathbb{Z}}
\newcommand{\Q}{\mathbb{Q}}
\newcommand{\R}{\mathbb{R}}
\newcommand{\C}{\mathbb{C}}
\newcommand{\F}{\mathbb{F}}
\newcommand{\A}{\mathbb{A}}
\newcommand{\PP}{\mathbb{P}}
\newcommand{\E}{\mathbb{E}}
\newcommand{\eps}{\varepsilon}
\newcommand{\dd}{\,\mathrm{d}}
\begin{document}
\section*{UnsolvedMath ID: 30001617}
\textbf{Source problem code:} catalogue entry
\subsection{Definitions and mathematical statement}
Let \(I_k(x)\) denote the nested dyadic intervals containing \(x\), with \(|I_k(x)|=2^{-k}\). For example, the standard half-open dyadic grid gives
\[
I_k(x)=\bigl[2^{-k}\lfloor 2^kx\rfloor,\,2^{-k}(\lfloor 2^kx\rfloor+1)\bigr).
\]
This specifies a conventional interpretation of the interval notation; the supplied record does not specify a grid or endpoint convention.

For a function \(f\) integrable on an interval \(I\), its interval mean is
\[
f_I=\frac{1}{|I|}\int_I f(t)\,\dd t.
\]
For real-valued \(f\) with these means defined, set
\[
d_k(x)=f_{I_{k+1}(x)}-f_{I_k(x)},\qquad
Q_n(x)=\sum_{k=n}^{\infty}d_k(x)^2\in[0,+\infty].
\]
Real-valuedness makes this series nonnegative. The source does not specify the scalar field or the admissible function class; real-valued local integrability is an interpretation sufficient to define these means, not a confirmed complete hypothesis of the intended problem.

The supplied question asks whether, for every function in the intended but unspecified admissible class,
\[
\limsup_{n\to\infty}
\frac{\left|f_{I_n(x)}-f(x)\right|}
{\sqrt{Q_n(x)\,\left|\log\left|\log Q_n(x)\right|\right|}}
<+\infty
\qquad\text{for Lebesgue-almost every }x\in\R.
\]
Here \(\log\) denotes the natural logarithm. Finiteness is pointwise almost everywhere; the displayed assertion does not require a common bound independent of \(x\) or \(f\).

\textbf{Unresolved qualifications.} The record supplies neither the function class nor conventions making the quotient meaningful at all relevant indices. In particular, the expression is undefined at \(Q_n(x)=0\), \(Q_n(x)=1\), or \(Q_n(x)=+\infty\), and its denominator vanishes at \(Q_n(x)=e^{-1}\) or \(Q_n(x)=e\). A condition such as \(0<Q_n(x)<e^{-1}\) eventually would make the displayed quotient well-defined, but is not supplied and is not imposed here. The case of eventually vanishing square-function tails requires an explicit convention or a separate formulation. Consequently, this is a statement template preserving the supplied question, pending recovery of its missing hypotheses and boundary conventions.
\subsection{Short English statement}
The question attributed in the supplied curation to Nicolau asks whether the error between a function and its dyadic interval means has a finite almost-everywhere limsup after normalization by its remaining dyadic square function and an iterated logarithm. The function class and conventions for degenerate square-function tails are missing, so the intended assertion requires clarification.
\subsection{Sources}
\begin{itemize}
\item Local UnsolvedMath registry, permanent ID \texttt{30001617}, source problem code \texttt{catalogue entry}, title: \emph{Law-of-Iterated-Logarithm Bounds for Dyadic Martingales}. The supplied record belongs to the Oberwolfach Reports Open Problems collection and gives proposed year 2010.
\item Supplied source citation: \emph{Operator Theory and Harmonic Analysis} (2011), Oberwolfach Reports, pp.\ 2813--2875, DOI \texttt{10.4171/owr/2010/49}. Source URL: \url{https://doi.org/10.4171/owr/2010/49}. The background describes Oberwolfach Reports as documenting workshops and their open-problem discussions.
\item Attribution: the supplied literature assessment describes this as the dyadic martingale square-function limsup question recorded by Nicolau. The proposer field is empty; the record does not establish that Nicolau proposed it.
\item Supplied literature source, labeled ``original OWR report'': \url{https://doi.org/10.4171/owr/2010/49}.
\item Literature assessment dated exactly \texttt{2026-08-21}: the supplied curation labels the problem open and reports related square-function estimates but no proof of this exact almost-everywhere assertion in the checked literature. This is supplied dated curation derived from the OWR record, not an independently verified literature assessment. The record itself requests independent verification of time-sensitive status and cited claims.
\item The supplied statement status is \texttt{corrected\_verified}, with a recorded manual semantic, typographic, and LaTeX review describing the cleanup as meaning-preserving and self-contained. These are supplied review labels; they do not resolve the missing hypotheses identified here. The report has not been independently consulted for this draft.
\end{itemize}
\end{document}
